\usepackage[usenames,dvipsnames]{xcolor}                                   

\usepackage{cancel}
\usepackage{amsfonts,amsmath,amssymb}
\usepackage[english]{babel}
\usepackage[normalem]{ulem}
\usepackage{hyperref}
\let\proof\relax 
\let\endproof\relax
\usepackage{amsthm}
\usepackage{graphicx}
\usepackage{xfrac}


\usepackage[linesnumbered,noend]{algorithm2e}

\usepackage{subcaption,caption}
\usepackage{epsfig}
\usepackage{ulem}
\usepackage{tikz}
\usepackage{mathrsfs}
\usepackage{bbold}             % This package was used to write the one with double bars 

\usepackage{soul}

\usepackage{pgfplots}

\pgfplotsset{compat = newest}

\usepackage{bm}
\usetikzlibrary{shapes,arrows}
\usetikzlibrary{matrix}
\usetikzlibrary{calc,patterns,angles,quotes,arrows,automata}
\graphicspath{ {./imgs/}}

\input{commands.tex}


\theoremstyle{theorem}
\newtheorem{lemma}{Lemma}
\newtheorem{proposition}{Proposition}
\newtheorem{conjecture}{Conjecture}
\newtheorem{theorem}{Theorem}
\newtheorem{corollary}{Corollary}

\theoremstyle{definition}
\newtheorem{definition}{Definition}
\newtheorem*{assumption*}{Assumptions}

\theoremstyle{remark}
\newtheorem{remark}{Remark}
\newtheorem{example}{Example}

\captionsetup{justification=Justified, format=plain, singlelinecheck=off}